\documentclass[11pt]{article}
\usepackage{amsmath, amssymb}
\usepackage{geometry}
\usepackage{hyperref}

% Load your ontology macros
\input{citation-ontology-mac-v2.0-infused.tex}

\begin{document}

% ============================================================
% TITLE PAGE
% ============================================================

\citeclass{FSFE-Title}
\provenanceanchor{FSFE-A6-TITLE-20261028}

\begin{center}
{\Large \textbf{Fun-Sphere Firewall Companion}}\\[6pt]
{\large A Reversible Physics Architecture for Geometric AI}\\[12pt]
Borealis S. Hedling\\
Dublin, Ireland\\
2026
\end{center}

\newpage

% ============================================================
% ABSTRACT
% ============================================================

\begin{citationmembrane}{A6-Abstract}
  \citeanchor{FSFEOverview}
  \citeanchor{ReversiblePhysics}
  \citeanchor{Holonomy}
  \citeanchor{Emb80}
\end{citationmembrane}

\section*{Abstract}

This work introduces the Fun-Sphere Firewall Companion (FSFE), a reversible
physics architecture for geometric AI. FSFE integrates reversible operator
algebra, holonomy-preserving geometric structures, an 80-dimensional embedding
manifold, altitude-layer documentation grammar, and a multi-altitude formal
verification stack. Runtime correspondence and completeness/soundness results
are outlined.

\newpage

% ============================================================
% INTRODUCTION
% ============================================================

\citeclass{FSFE-Intro}
\provenanceanchor{FSFE-A6-INTRO-20261028}

\section{Introduction}

\subsection{Motivation}
Reversible physics provides structural guarantees for geometric AI systems.

\subsection{FSFE Overview}
FSFE consists of:
\begin{itemize}
  \item $\Sigma/\Lambda$ reversible counters
  \item Quest Algebra branching
  \item Tensor Coupling
  \item Branch Stability
  \item Holonomy geometry
  \item 80D embedding manifold
  \item Altitude documentation system
  \item Tile grammar
\end{itemize}

\subsection{Contributions}
A concise list of FSFE contributions.

\newpage

% ============================================================
% SYSTEM ARCHITECTURE
% ============================================================

\begin{citationmembrane}{A6-SystemArchitecture}
  \citeanchor{AltitudeMap}
  \citeanchor{OperatorSuite}
  \citeanchor{HolonomyModel}
  \citeanchor{EmbeddingManifold}
\end{citationmembrane}

\section{System Architecture Overview}

\subsection{Altitude Layers}
Conceptual separation A1--A12.

\subsection{Reversible Physics Engine}
Operator suite and state machine.

\subsection{Geometric Structures}
Holonomy and embedding manifold.

\subsection{Documentation Grammar}
Tile system and altitude rules.

\newpage

% ============================================================
% FORMAL DEFINITIONS
% ============================================================

\citeclass{FSFE-Definitions}
\provenanceanchor{FSFE-A6-DEFINITIONS-20261028}

\section{Formal Definitions}

\subsection{FSFEState}
\begin{verbatim}
FSFEState := {
  sigma      : Nat,
  lambda     : Nat,
  qa         : QAState,
  tensor     : TensorState,
  branch     : BranchState,
  holonomy   : HolonomyState,
  embedding  : Emb80
}
\end{verbatim}

\subsection{Reversible Operators}
\begin{verbatim}
Σ_inc(s) := { s with sigma := s.sigma + 1 }
Λ_inc(s) := { s with lambda := s.lambda + 1 }
QuestStep(s) := { s with qa := QA_next(s.qa) }
TensorCouple(s) := { s with tensor := Tensor_couple(s.tensor) }
StabilizeBranch(s) := { s with branch := Branch_stable(s.branch) }
MobileTransition(s, t) := Transition_apply(s, t)
\end{verbatim}

\subsection{Invariants}
\begin{verbatim}
Inv_struct(s) := ...
Inv_hol(s) := ...
Inv_emb(s) := ...
FSFEInvariant(s) := ...
\end{verbatim}

\subsection{Holonomy Model}
\begin{verbatim}
HolonomyState := {
  connection : Conn,
  curvature  : Curv,
  transport  : Transport,
  group      : HolGroup
}
\end{verbatim}

\subsection{Embedding Manifold}
\begin{verbatim}
Emb80 := {
  coords : Fin 80 → Real,
  metric : Metric80,
  chart  : Chart80,
  trans  : Transition80
}
\end{verbatim}

\subsection{Altitude Semantics}
\begin{verbatim}
Altitude := { A1, A2, ..., A12 }
AltitudeCompatible(a_high, a_low) :=
  altitude_index(a_low) ≤ altitude_index(a_high)
\end{verbatim}

\subsection{Tile Grammar}
\begin{verbatim}
Tile :=
  | Concept | Operator | Invariant | Transition
  | Algebra | Geometry | Runtime | Telemetry | Meta
\end{verbatim}

\subsection{Documentation Composition}
\begin{verbatim}
Document := { tiles : List Tile }
DocumentValid(d) := ...
\end{verbatim}

\newpage

% ============================================================
% THEOREM STATEMENTS
% ============================================================

\begin{citationmembrane}{A6-Theorems}
  \citeanchor{InvariantTheorems}
  \citeanchor{OperatorAlgebraTheorems}
  \citeanchor{HolonomyTheorems}
  \citeanchor{EmbeddingTheorems}
\end{citationmembrane}

\section{Theorem Statements}

\subsection{State Invariant Theorems}
\begin{verbatim}
Theorem StructuralInvariantPreserved : ...
Theorem HolonomyInvariantPreserved : ...
Theorem EmbeddingInvariantPreserved : ...
\end{verbatim}

\subsection{Operator Algebra Theorems}
\begin{verbatim}
Theorem OperatorComposeAssociative : ...
Theorem OperatorIdentityNeutral : ...
Theorem SigmaLambdaReversible : ...
\end{verbatim}

\subsection{Quest Algebra Theorems}
\begin{verbatim}
Theorem QuestStepProgress : ...
Theorem QuestTensorStability : ...
\end{verbatim}

\subsection{Tensor Coupling Theorems}
\begin{verbatim}
Theorem TensorRankPreserved : ...
Theorem TensorDimensionStable : ...
\end{verbatim}

\subsection{Branch Stability Theorems}
\begin{verbatim}
Theorem BranchStabilization : ...
Theorem BranchMobileStable : ...
\end{verbatim}

\subsection{Holonomy Theorems}
\begin{verbatim}
Theorem HolonomyGroupClosed : ...
Theorem ParallelTransportCompatible : ...
\end{verbatim}

\subsection{Embedding Manifold Theorems}
\begin{verbatim}
Theorem ChartTransitionValid : ...
Theorem MetricPreserved : ...
\end{verbatim}

\subsection{Altitude Semantics Theorems}
\begin{verbatim}
Theorem AltitudeCompatReflexive : ...
Theorem AltitudeCompatTransitive : ...
\end{verbatim}

\subsection{Tile Grammar Theorems}
\begin{verbatim}
Theorem TileSequenceStable : ...
Theorem MetaTileCompatible : ...
\end{verbatim}

\subsection{Documentation Composition Theorems}
\begin{verbatim}
Theorem DocumentCompositionStable : ...
Theorem DocumentAltitudePreserved : ...
\end{verbatim}

\newpage

% ============================================================
% RUNTIME CORRESPONDENCE
% ============================================================

\citeclass{FSFE-Runtime}
\provenanceanchor{FSFE-A6-RUNTIME-20261028}

\section{Runtime Correspondence}

\subsection{Python/Rust Reversible Runtime}
\subsection{Telemetry Adapter}
\subsection{Constellation Engine}
\subsection{Reversible State Projection}
\subsection{Runtime Invariant Preservation}

\newpage

% ============================================================
% GEOMETRY
% ============================================================

\begin{citationmembrane}{A6-Geometry}
  \citeanchor{Holonomy}
  \citeanchor{Curvature}
  \citeanchor{ParallelTransport}
\end{citationmembrane}

\section{Geometry}

\subsection{Holonomy}
\subsection{Curvature}
\subsection{Parallel Transport}
\subsection{Holonomy Group}
\subsection{Holonomy Invariants}

\newpage

% ============================================================
% EMBEDDING MANIFOLD
% ============================================================

\citeclass{FSFE-Embedding}
\provenanceanchor{FSFE-A6-EMBEDDING-20261028}

\section{Embedding Manifold}

\subsection{Coordinate Charts}
\subsection{Metric}
\subsection{Transition Maps}
\subsection{Embedding Invariants}
\subsection{Embedding Preservation}

\newpage

% ============================================================
% VERIFICATION
% ============================================================

\begin{citationmembrane}{A6-Verification}
  \citeanchor{FSFEAxioms}
  \citeanchor{Completeness}
  \citeanchor{Soundness}
  \citeanchor{Normalization}
  \citeanchor{AltitudeTransitions}
\end{citationmembrane}

\section{Verification}

\subsection{Axioms (A5)}
\subsection{Completeness (A5)}
\subsection{Soundness (A7--A12)}
\subsection{Normalization}
\subsection{Altitude Transitions}
\subsection{Tile Composition}

\newpage

% ============================================================
% APPENDIX STRUCTURE
% ============================================================

\citeclass{FSFE-Appendix}
\provenanceanchor{FSFE-A6-APPENDIX-20261028}

\section{Appendix Structure}

\begin{verbatim}
Appendix A — Operator Algebra Expansion
Appendix B — Holonomy Tensor Derivations
Appendix C — Embedding Manifold Charts
Appendix D — Transition Membrane Formalism
Appendix E — Telemetry Projection Math
Appendix F — Altitude Grammar Formalization
Appendix G — Runtime Correspondence Proofs
\end{verbatim}

\end{document}
